%% ma: L12-B14-0036
%% lop: 12
%% chuong: 5
%% bai: 14
%% dang: 12.14.19
%% loai: TN4
%% muc: TH
%% dap_an: "D"
%% nguon: "CĐ2.5.pdf giáo viên gửi 27/10/2026 (M07, Lê Minh Khôi – Nguyễn Anh Khôi, Chương II Chủ đề 5 Tích có hướng của hai vectơ), Câu 20"
%% kien_thuc: "Bài 14 (tích có hướng, phương trình mặt phẳng, khoảng cách từ điểm đến mặt phẳng), lớp 11 (thể tích khối chóp, khối hộp)"
%% trang_thai: chinh-thuc
%% ghi_chu: "Giáo viên đã duyệt 27/10/2026"
\begin{ex}
Trong không gian $Oxyz$, cho tứ diện $ABCD$ có $A(1;-2;0)$, $B(2;0;3)$, $C(-2;1;3)$ và $D(0;1;1)$. Thể tích của tứ diện $ABCD$ bằng
\choice{$6$}{$8$}{$12$}{\True $4$}
\loigiai{$\vec n=\left[\overrightarrow{AB},\overrightarrow{AC}\right]=(-3;-12;9)$, $\vec n\cdot\overrightarrow{AD}=-24$.

$V=\dfrac13\cdot\dfrac12|\vec n|\cdot\dfrac{24}{|\vec n|}=4$.}
\end{ex}
